\section{Walking out of the Forest}

\begin{quotex}
Flectere si nequeo superos, Acheronta movebo.
\flright{\textsc{Virgil}}

formans lucem et creans tenebras, faciens pacem et creans malum: ego Dominus faciens omnia haec.
\flright{\textsc{Isaias 45:7}}
\end{quotex}


When Dante found himself lost in a forest, he had to rely of a pagan, Virgil, to lead him out. Specifically, it was not a catechism, not a prayer, i.e., not anything supernatural that he relied on. Virgil led him through the most terrifying part of the journey, through Hell.

Virgil resides in Limbo, where he would experience the natural happiness of the pagans. However, the mission of Christ was to restore the human state to that before the Fall, which is supernatural happiness. With roots more than 10 centuries before the founding of Rome, there was some tension in grafting that older Tradition onto the newer ones of Rome, not to mention the rest of Europe.

Now that the essence of Christianity has been lost, it appears alien. Even soi-disant ``conservative Catholics'' have no experience of Medieval Catholicism and would probably reject much of its worldview.

\paragraph{Heaven and Earth}


\begin{quotex}
In the beginning God created the heaven and the earth.
\flright{\textsc{Moses}, \emph{Genesis 1:1}}

The nameless is the origin of Heaven and Earth.
\flright{\textsc{Lao Tzu}, \emph{Tao Te Ching 1:3}}
\end{quotex}


Moses (1393 BC -- 1273 BC) and Lao Tzu (\textasciitilde{} 500 BC) seem to be onto something here. The same insight from vastly different space and time shows that is it not a vague dogma to be believed, but rather a deep truth. Nevertheless, it is often poorly understood, so it should be meditated on, for as long as it takes. It holds the secret behind vertical thinking (as opposed to horizontal). Its analogues are essence/substance, idea/form, ideal/real, Purusha/Prakriti, thinking/sensing, the worlds of Being and Becoming.

\paragraph{Being and Beyond Being}


\begin{quotex}
Heaven and earth do not act from any wish to be benevolent; they deal with all things as straw dogs are dealt with. The sages do not act from any wish to be benevolent.
\flright{\textsc{Tao Te Ching}}
\end{quotex}


This calls for a decision. According to \textbf{Rene Guenon}, Being is the first affirmation of Beyond Being (or Non-Being), and any moral considerations are out of place (as described by Lao Tzu). That seems to be what Moses revealed: God is Being. Yet, Plato claimed that the Good is ultimate, not Being. And St. John takes a stronger position that God is Love. \textbf{Valentin Tomberg} raises this question:

\begin{quotex}
Yet what is the full significance of the adoption of the primacy of being, instead of that of good, or according to St. John, that of love? The idea of being is neutral from the point of view of the moral life. There is no need to have the experience of the good and the beautiful in order to arrive at it. The experience solely of the mineral realm already suffices to arrive at the morally neutral idea of being.
\end{quotex}


The universe of Being is the tour de force of physics and mathematics. However, it is a universe without love, without life, without consciousness, and without moral life. Those are all mysteries without explanation for science. Nevertheless, the scientist himself experiences all those things.

The universe of Being suits minerals but omits the human. Hence, it is more appropriate to say that Being is the first affirmation of Love, or the Good.

\paragraph{Infinite Possibilities}


\begin{quotex}
It is of the essence of spirit to embody itself.
\flright{\textsc{Robin Collingwood}}
\end{quotex}


The realm of ideas subsists in the mind of God. This accounts for the Infinity of God, since there is no limit. There are two kinds of ideas:

\begin{itemize}


\item Possibilities of non-manifestation

\item Possibilities of manifestation
\end{itemize}


The possibilities of manifestation are actualized in existence, at some level (not just the physical level). Obviously, no all possibilities can manifest at the same place and time; they need to be compossible with each other. That is why events are restricted to definite times and space.

Now, a possibility of manifestation must manifest, otherwise the idea makes no sense.

The possibilities of non-manifestation can never be actualized in existence. For example, there are logical contradictions; also, the Void and Silence cannot be manifested because they are the opposite of being.

\paragraph{Free Will and Necessity}


Intelligence is the ability to make choose one's ends and the means to achieve them. That is practical reason and is a requirement in the world of Love. This not true of the world of Being; robots, for example, are minerals and act out their prior programming.

Free will is not the power to act arbitrarily, as, for example, in the manner of a Sartrean existentialist. In one of his novels, a character drove a knife into his hand just to demonstrate his freedom.

Hence, free will is the power to act from one's essence. Its opposite is compulsion from outside oneself. As a man develops his essence, he becomes purer, and his will becomes more free.

So free will is not simply the power to make choices. \emph{A fortiori}, this is true of God. To imagine that God is eternally making choices is to attribute to him a gnomic will.

\paragraph{Evil and Privation}


The idea the evil is not a positive quality is confusing to many people who might bother to think about it. They think it means that evil is imaginary or its effects are not real. Rather, it means that evil is the absence of goodness.

Compare, for example, the notion of darkness, which is the absence of light. If you carry a torch to try to find the dark, you will never succeed. Before you go to sleep at night, you don't turn on the darkness, instead you turn off the light.

Evil is what remains when you don't choose the Good, whether out of weakness, ignorance, malice, or cupidity. It follows, then, that if everyone were Good, there would be no evil. Hence, evil has no being in itself.

\paragraph{Natural Evil}


It is difficult to understand what this could even mean. Nature is the world understood by science, i.e., the world of Being. Scientists seek to understand its laws and causes and moral considerations are never part of their equations.

If a giant meteor crashes into the Moon, leaving a crater, no one describes that as a ``natural evil''. However, if it crashes into Tokyo, killing hundreds of thousands or people, then the description might change.

There is no agent in natural phenomena, so there can be no question of morality. Rather, it is up to science, technology, and engineering to learn how to deal with disastrous natural events. For example, better predictive models, sounder building techniques, medical advances, and so on, are the appropriate responses.

\paragraph{Pain, Suffering and Hell}


For the secularist, pain and suffering are the worst things that can befall a human being. Not to diminish the horror of intense suffering, we affirm instead that the pains of Hell are much worse. Transient human suffering pales in relationship to the loss of one's soul.

\paragraph{Appendix}


\textbf{St. Ambrose: Concerning repentance (Number 27)}


What the Apostle means by the rod is shown by his invective against fornication, his denunciation of incest, his reprehension of pride, because they were puffed up who ought rather to be mourning, and lastly, his sentence on the guilty person, that he should be excluded from communion, and delivered to the adversary, not for the destruction of the soul but of the flesh. For as the Lord did not give power to Satan over the soul of holy Job, but allowed him to afflict his body, so here, too, the sinner is delivered to Satan for the destruction of the flesh, that the serpent might lick the dust of his flesh, but not hurt his soul. Let, then, our flesh die to lusts, let it be captive, let it be subdued, and not war against the law of our mind, but die in subjection to a good service, as in Paul, who buffeted his body that he might bring it into subjection, in order that his preaching might become more approved, if the law of his flesh agreed and was consonant with the law of his flesh. For the flesh dies when its wisdom passes over into the spirit, so that it no longer has a taste for the things of the flesh, but for the things of the spirit. Would that I might see my flesh growing weak, would that I were not dragged captive into the law of sin, would that I lived not in the flesh, but in the faith of Christ! And so there is greater grace in the infirmity of the body than in its soundness.


\flrightit{Posted on 2020-04-05 by Cologero}

